% Resumen del mundo «Fracciones Algebraicas» (Álgebra, mundo 9: expresiones racionales y fracciones parciales).
% Lecciones alg-25a, alg-25b, alg-101, alg-102. Ejemplos de los videos (el sendero de Beto, el estanque de Ana y Beto,
% el bote de Caro y la curva de dos ríos); el taller trae ejercicios nuevos (no repite los de la app, del jefe ni de la evaluación).
\documentclass[11pt,a4paper]{article}
\usepackage{zuka-resumen}
\zukacurso{algebra}
\zukamundo{9}{Fracciones Algebraicas}{Expresiones racionales y fracciones parciales}

\newcommand{\gr}{\operatorname{gr}}
\newcommand{\Dom}{\operatorname{Dom}}

\begin{document}

\portada{%
Una \clave{expresión racional} es una fracción de polinomios. Casi siempre se comporta bien, pero donde el
denominador se anula aparece un \clave{río}: la curva se dispara junto a una asíntota. En este mundo aprendiste a
decidir dónde existe una fracción algebraica, a simplificarla y operarla, a resolver sistemas de ecuaciones y a
partir una fracción en fracciones parciales. Este resumen reúne las definiciones y los ejemplos de los videos, y
termina con un taller de práctica.}{%
\item \textbf{Nivel 9-1.} Dominio y simplificación: denominador distinto de cero, cancelar factores, asíntotas y huecos.
\item \textbf{Nivel 9-2.} Operaciones con fracciones algebraicas: producto, cociente, suma con denominador común y fracciones complejas.
\item \textbf{Nivel 9-3.} Sistemas de ecuaciones lineales: interpretación gráfica, reducción, sustitución, igualación y escalonamiento.
\item \textbf{Nivel 9-4.} Sistemas no lineales y fracciones parciales.}

\begin{clavebox}[LAS TRES IDEAS QUE SOSTIENEN TODO EL MUNDO]
\renewcommand{\arraystretch}{1.5}\begin{tabular}{@{}c@{\qquad}c@{\qquad}c@{}}
$\dfrac{P(x)}{Q(x)}$ con $Q(x)\neq0$ & $\dfrac{A\,C}{B\,C}=\dfrac{A}{B}$ & solución $=$ punto de corte\\[2pt]
\small nunca se divide entre cero & \small se cancela lo que multiplica & \small dos condiciones a la vez
\end{tabular}
\end{clavebox}

%=====================================================================
\nivel{9-1}{Dominio y simplificación}

Beto camina por un sendero junto al río. El sendero sigue una fracción de polinomios,
\[ f(x)=\frac{x^2-1}{x^2+2x-3}, \]
pero en dos puntos la fórmula falla.

\begin{definicion}[Expresión racional y dominio]
Una \clave{expresión racional} es un cociente $\dfrac{P(x)}{Q(x)}$ de polinomios, con $Q$ distinto del polinomio cero.
Su \clave{dominio} es el conjunto de los reales donde el denominador no se anula:
\[ \Dom=\conjunto{x\in\R \tq Q(x)\neq 0}. \]
\end{definicion}
\textbf{Por qué no se divide entre cero.} $6\div 2=3$ porque $2\cdot3=6$. Para $6\div0$ habría que hallar un número
que multiplicado por $0$ dé $6$, y no existe ninguno.

\begin{propiedad}[Simplificar]
Si $C(x)\neq0$, entonces $\dfrac{A(x)\,C(x)}{B(x)\,C(x)}=\dfrac{A(x)}{B(x)}$. Se factoriza arriba y abajo y se
cancelan \emph{factores} comunes. El dominio se lee siempre en la expresión \emph{original}.
\end{propiedad}

\begin{ejemplo}[El sendero de Beto]
Se factoriza arriba y abajo:
\[ \frac{x^2-1}{x^2+2x-3}=\frac{(x-1)(x+1)}{(x-1)(x+3)}=\frac{x+1}{x+3},\qquad x\neq1,\ x\neq-3 . \]
El denominador original se anula en $x=1$ y en $x=-3$, así que $\Dom=\R\setminus\{-3,\,1\}$. Pero los dos puntos
no fallan igual:
\begin{itemize}
\item En $x=-3$ solo se anula el denominador: la expresión se dispara. Hay una \clave{asíntota vertical} (el río).
\item En $x=1$ se anulan los dos y el factor $x-1$ se cancela: la gráfica es la de $\frac{x+1}{x+3}$ con un
      \clave{hueco} en $\left(1,\tfrac{2}{4}\right)=\left(1,\tfrac12\right)$.
\end{itemize}
\end{ejemplo}

\begin{minipage}[c]{0.48\linewidth}
\begin{definicion}[Asíntotas]
La recta $x=a$ es \clave{asíntota vertical} si $\abs{f(x)}$ crece sin límite cuando $x$ se acerca a $a$; en una
fracción simplificada, ocurre donde se anula el denominador.\par\smallskip
La recta $y=L$ es \clave{asíntota horizontal} si $f(x)$ se acerca a $L$ cuando $\abs{x}$ crece sin límite.
\end{definicion}
\end{minipage}\hfill
\begin{minipage}[c]{0.48\linewidth}\centering
\begin{tikzpicture}
\begin{axis}[zplot,width=7.4cm,height=5.8cm,xmin=-7,xmax=4,ymin=-6,ymax=6,
  xtick={-6,-3,3},ytick={-4,1,4},xlabel={$x$},ylabel={$y$},
  xlabel style={at={(ticklabel* cs:1)},anchor=north west},ylabel style={at={(ticklabel* cs:1)},anchor=south west}]
  \addplot[zcoral,dashed,line width=0.9pt] coordinates {(-3,-6) (-3,6)};
  \addplot[zmute,dashed,line width=0.9pt] coordinates {(-7,1) (4,1)};
  \addplot[zaccent,domain=-7:-3.4,samples=80] {(x+1)/(x+3)};
  \addplot[zaccent,domain=-2.714:4,samples=120] {(x+1)/(x+3)};
  \addplot[only marks,mark=*,mark size=2.4pt,fill=white,draw=zaccent,line width=1pt] coordinates {(1,0.5)};
  \node[font=\footnotesize,text=zcoral,anchor=west] at (axis cs:-2.9,-5.2) {$x=-3$};
  \node[font=\footnotesize,text=zmute,anchor=north] at (axis cs:-5.5,0.85) {$y=1$};
  \draw[zmute,line width=0.6pt] (axis cs:1.15,0.3) -- (axis cs:2.2,-1.8) node[below,font=\footnotesize,text=zink] {hueco};
\end{axis}
\end{tikzpicture}

{\small\color{zmute}$y=\frac{x^2-1}{x^2+2x-3}$: un río en $x=-3$ y un hueco en $x=1$.}
\end{minipage}

\begin{propiedad}[Asíntota horizontal de $\frac{P}{Q}$, con $\gr P=n$ y $\gr Q=m$]
Si $n<m$, la asíntota es $y=0$. Si $n=m$, es $y=\dfrac{\text{coef.\ principal de }P}{\text{coef.\ principal de }Q}$.
Si $n>m$, no hay asíntota horizontal.
\end{propiedad}
En el sendero, $n=m=2$ y la asíntota es $y=\frac11=1$: para $x$ enorme, los números sueltos dejan de pesar.

\begin{ojo}[Solo se cancelan factores]
$\dfrac{x+5}{5}\neq x$: el $5$ de arriba está \emph{sumando}, no multiplica a todo el numerador. Y
$\dfrac{x^2+3}{x^2}\neq 3$ por la misma razón. Antes de cancelar, factoriza.
\end{ojo}

%=====================================================================
\nivel{9-2}{Operaciones con fracciones algebraicas}

Ana sola llena el estanque en $3$ horas; Beto solo, en $6$. ¿Cuánto tardan juntos? No es el promedio: con ayuda,
Ana tardaría \emph{menos} de $3$ horas.

\begin{ejemplo}[El estanque, con números]
Se piensa en cuánto llena cada uno \emph{por hora}: Ana $\frac13$ de estanque y Beto $\frac16$. Con el denominador
común $6$:
\[ \frac13+\frac16=\frac26+\frac16=\frac36=\frac12\ \text{ de estanque por hora}. \]
Para llenar un estanque entero: $1\div\frac12=2$ horas.
\end{ejemplo}

\begin{figura}[Por hora, Ana llena dos sextos del estanque y Beto uno: juntos, la mitad.]
\begin{tikzpicture}[x=1.4cm,y=1cm]
  \fill[zsky!45] (0,0) rectangle (2,0.8);
  \fill[zsun!55] (2,0) rectangle (3,0.8);
  \foreach \i in {0,...,6} \draw[zink,line width=0.6pt] (\i,0) -- (\i,0.8);
  \draw[zink,line width=0.9pt] (0,0) rectangle (6,0.8);
  \draw[decorate,decoration={brace,amplitude=5pt,mirror},zmute] (0,-0.1) -- (2,-0.1) node[midway,below=5pt,font=\footnotesize,text=zink] {Ana: $\tfrac13=\tfrac26$};
  \draw[decorate,decoration={brace,amplitude=5pt,mirror},zmute] (2,-0.1) -- (3,-0.1) node[midway,below=5pt,font=\footnotesize,text=zink] {Beto: $\tfrac16$};
  \draw[decorate,decoration={brace,amplitude=5pt},zmute] (0,0.9) -- (3,0.9) node[midway,above=5pt,font=\footnotesize,text=zink] {juntos: $\tfrac36=\tfrac12$};
  \node[font=\footnotesize,text=zmute,anchor=west] at (6.1,0.4) {un estanque};
\end{tikzpicture}
\end{figura}

\begin{propiedad}[Las cuatro operaciones]
Para expresiones con denominadores no nulos:
\[
\frac{A}{B}\cdot\frac{C}{D}=\frac{AC}{BD},\qquad
\frac{A}{B}\div\frac{C}{D}=\frac{A}{B}\cdot\frac{D}{C}\ (C\neq0),\qquad
\frac{A}{B}\pm\frac{C}{B}=\frac{A\pm C}{B}.
\]
Para sumar o restar con denominadores distintos se usa el \clave{mínimo común denominador} (mcd): se factoriza cada
denominador y se toma cada factor una vez, con su mayor exponente; luego cada fracción se amplifica por lo que le falta.
\end{propiedad}

\begin{ejemplo}[El estanque, con letras]
Si Ana tarda $x$ horas y Beto $3$ horas más, por hora llenan $\frac1x$ y $\frac1{x+3}$. El mcd es $x(x+3)$:
\[ \frac1x+\frac1{x+3}=\frac{x+3}{x(x+3)}+\frac{x}{x(x+3)}=\frac{2x+3}{x(x+3)} . \]
El tiempo juntos es $1$ entre esa suma: se voltea la fracción,
\[ T(x)=\frac{x(x+3)}{2x+3}. \qquad\text{Con }x=3:\quad T=\frac{3\cdot6}{9}=2\text{ horas, como con números.} \]
\end{ejemplo}

\begin{ejemplo}[Una fracción compleja]
Se resuelve arriba y abajo por separado y luego se divide:
\[ \frac{1+\frac1x}{1-\frac1{x^2}}=\frac{\frac{x+1}{x}}{\frac{x^2-1}{x^2}}=\frac{x+1}{x}\cdot\frac{x^2}{(x-1)(x+1)}=\frac{x}{x-1},
\qquad x\neq0,\ \pm1 . \]
\end{ejemplo}

\begin{ojo}[Errores al operar]
\begin{itemize}
\item $\dfrac{a}{b}+\dfrac{c}{d}\neq\dfrac{a+c}{b+d}$: con $\frac12+\frac12$ saldría $\frac24$, ¡menos que uno!
\item Al restar, el signo afecta a \emph{todo} el numerador: $\dfrac{A}{B}-\dfrac{C+D}{B}=\dfrac{A-C-D}{B}$.
\item $2-x$ y $x-2$ no son iguales sino opuestos: $2-x=-(x-2)$.
\end{itemize}
\end{ojo}

%=====================================================================
\nivel{9-3}{Sistemas de ecuaciones lineales}

Caro navega por el río. A favor de la corriente va a $\SI{15}{km/h}$; en contra, solo a $\SI{9}{km/h}$. Si $v$ es la
velocidad del bote y $c$ la de la corriente:
\[ \begin{cases} v+c=15\\ v-c=9 \end{cases} \]

\begin{minipage}[c]{0.5\linewidth}
Con una sola pista no alcanza: $v+c=15$ tiene muchas soluciones ($10$ y $5$, $12$ y $3$, $14$ y $1$\dots), y todas
forman una \clave{recta}. La segunda pista es otra recta. Solo el punto donde se \clave{cortan} cumple las dos:
$v=12$, $c=3$.
\begin{definicion}[Solución de un sistema]
Una \clave{solución} de un sistema es una asignación de valores a las incógnitas que cumple \emph{todas} las
ecuaciones a la vez.
\end{definicion}
\end{minipage}\hfill
\begin{minipage}[c]{0.46\linewidth}\centering
\begin{tikzpicture}
\begin{axis}[zplot,width=7cm,height=5.6cm,xmin=0,xmax=18,ymin=-2,ymax=10,
  xtick={3,6,9,12,15},ytick={3,6,9},xlabel={$v$},ylabel={$c$},
  xlabel style={at={(ticklabel* cs:1)},anchor=north west},ylabel style={at={(ticklabel* cs:1)},anchor=south west}]
  \addplot[zsky,domain=5:17.5] {15-x};
  \addplot[zcoral,domain=7.5:18] {x-9};
  \addplot[only marks,mark=*,mark size=2.4pt,zink] coordinates {(12,3)};
  \node[font=\footnotesize,text=zsky,anchor=west] at (axis cs:7.6,8) {$v+c=15$};
  \node[font=\footnotesize,text=zcoral,anchor=east] at (axis cs:14.8,6.5) {$v-c=9$};
  \node[font=\footnotesize,text=zink,anchor=west] at (axis cs:12.4,3.4) {$(12,3)$};
\end{axis}
\end{tikzpicture}
\end{minipage}

\subsection{Tres caminos sin dibujar}
\begin{center}
\begin{tabularx}{\linewidth}{@{}>{\bfseries}l L L@{}}
\toprule
Método & Idea & Con el bote de Caro\\ \midrule
Reducción & sumar o restar ecuaciones para que una incógnita desaparezca & sumando: $2v=24$, $v=12$; luego $c=3$\\
Sustitución & despejar una incógnita y reemplazarla en la otra ecuación & $v=15-c$ en la segunda: $15-2c=9$, $c=3$\\
Igualación & despejar la misma incógnita en las dos e igualar & $15-c=9+c$, así que $c=3$\\
\bottomrule
\end{tabularx}
\end{center}

\subsection{¿Cuántas soluciones?}
\begin{figura}[Dos rectas del plano: se cortan, son paralelas o son la misma.]
\begin{tikzpicture}[scale=0.8]
  \foreach \x/\t/\s in {0/{una solución}/{se cortan},5.5/{ninguna}/{paralelas distintas},11/{infinitas}/{la misma recta}} {
    \draw[zline] (\x-2,-1.4) rectangle (\x+2,1.4);
    \node[font=\small\bfseries] at (\x,-1.9) {\t};
    \node[font=\footnotesize,text=zmute] at (\x,-2.45) {\s};}
  \draw[zsky,line width=1.2pt] (-1.8,-1.1) -- (1.8,1.1);
  \draw[zcoral,line width=1.2pt] (-1.8,1.0) -- (1.8,-1.0);
  \fill[zink] (0,0) circle (2.2pt);
  \draw[zsky,line width=1.2pt] (3.7,-1.2) -- (6.6,1.2);
  \draw[zcoral,line width=1.2pt] (4.4,-1.2) -- (7.3,1.2);
  \draw[zsky,line width=2.6pt] (9.2,-1.1) -- (12.8,1.1);
  \draw[zcoral,line width=1.1pt,dashed] (9.2,-1.1) -- (12.8,1.1);
\end{tikzpicture}
\end{figura}
En $a_1x+b_1y=c_1$, $a_2x+b_2y=c_2$: si $\frac{a_1}{a_2}=\frac{b_1}{b_2}$, las rectas son paralelas; no hay solución
si además $\frac{c_1}{c_2}$ es distinto, y hay infinitas si también coincide (una ecuación es múltiplo de la otra).

\subsection{Tres incógnitas: la escalera}
\begin{ejemplo}[Escalonar un sistema $3\times3$]
\[
\begin{cases} x+y+z=9\\ 2x-y+z=7\\ x+2y-z=8 \end{cases}
\xrightarrow[\;E_3-E_1\;]{\;E_2-2E_1\;}
\begin{cases} x+y+z=9\\ -3y-z=-11\\ y-2z=-1 \end{cases}
\xrightarrow{\;E_2+3E_3\;}
\begin{cases} x+y+z=9\\ y-2z=-1\\ -7z=-14 \end{cases}
\]
Abajo queda $z=2$. Un escalón arriba, $y=-1+2z=3$. Arriba, $x=9-3-2=4$. Comprobación en la segunda ecuación:
$2(4)-3+2=7$.
\end{ejemplo}

\begin{ojo}[Comprueba siempre]
Reemplaza la solución en \emph{todas} las ecuaciones originales, no solo en la última que usaste. Y si al reducir
aparece $0=5$, el sistema no tiene solución; si aparece $0=0$, tiene infinitas.
\end{ojo}

%=====================================================================
\nivel{9-4}{Sistemas no lineales y fracciones parciales}

\subsection{Sistemas no lineales}
\begin{minipage}[c]{0.52\linewidth}
¿Dónde se cortan una orilla curva, $y=x^2$, y un puente recto, $y=2x+3$? Donde las dos alturas son iguales:
\begin{align*}
x^2&=2x+3\\
x^2-2x-3&=0\\
(x-3)(x+1)&=0 .
\end{align*}
Con $x=3$, $y=9$; con $x=-1$, $y=1$. Los cortes son $(3,9)$ y $(-1,1)$. En general se \clave{despeja una incógnita}
de la ecuación más sencilla y se \clave{sustituye} en la otra; el discriminante de la cuadrática que resulta dice si
hay dos, uno o ningún corte.
\end{minipage}\hfill
\begin{minipage}[c]{0.44\linewidth}\centering
\begin{tikzpicture}
\begin{axis}[zplot,width=6.6cm,height=6cm,xmin=-2.6,xmax=4,ymin=-2,ymax=12,
  xtick={1,3},ytick={1,5,9},xlabel={$x$},ylabel={$y$},
  xlabel style={at={(ticklabel* cs:1)},anchor=north west},ylabel style={at={(ticklabel* cs:1)},anchor=south west}]
  \addplot[zaccent,domain=-2.5:3.45,samples=100] {x^2};
  \addplot[zcoral,domain=-2.5:4] {2*x+3};
  \addplot[only marks,mark=*,mark size=2.4pt,zink] coordinates {(3,9) (-1,1)};
  \node[font=\footnotesize,text=zink,anchor=east] at (axis cs:2.85,9.6) {$(3,9)$};
  \node[font=\footnotesize,text=zink,anchor=north] at (axis cs:-1,-0.4) {$(-1,1)$};
  \node[font=\footnotesize,text=zaccent,anchor=west] at (axis cs:-2.4,10.8) {$y=x^2$};
  \node[font=\footnotesize,text=zcoral,anchor=west] at (axis cs:0.2,11.4) {$y=2x+3$};
\end{axis}
\end{tikzpicture}
\end{minipage}

\subsection{Fracciones parciales}
Beto dibuja el mapa del río: una curva enredada entre dos ríos, en $x=1$ y $x=-1$,
\[ g(x)=\frac{3x+1}{x^2-1}. \]
Esa curva esconde un secreto: es la suma de dos fracciones sencillas. Partirla es el camino inverso de sumar.

\begin{definicion}[Descomposición en fracciones parciales]
Si $\frac{P}{Q}$ es \clave{propia} ($\gr P<\gr Q$), se factoriza $Q$ en factores lineales y cuadráticos irreducibles,
y cada factor aporta fracciones según esta tabla:
\begin{center}
\begin{tabular}{@{}ll@{}}
\toprule
Factor de $Q$ & Fracciones que aporta\\ \midrule
$(x-a)$, lineal & $\dfrac{A}{x-a}$\\[8pt]
$(x-a)^k$, lineal repetido & $\dfrac{A_1}{x-a}+\dfrac{A_2}{(x-a)^2}+\dots+\dfrac{A_k}{(x-a)^k}$\\[8pt]
$ax^2+bx+c$, irreducible ($b^2-4ac<0$) & $\dfrac{Bx+C}{ax^2+bx+c}$\\[6pt]
\bottomrule
\end{tabular}
\end{center}
\end{definicion}

\begin{ejemplo}[La curva de dos ríos]
Como $x^2-1=(x-1)(x+1)$, se busca
\[ \frac{3x+1}{(x-1)(x+1)}=\frac{A}{x-1}+\frac{B}{x+1}\quad\Longrightarrow\quad 3x+1=A(x+1)+B(x-1). \]
\textbf{Comparando coeficientes:} $3x+1=(A+B)x+(A-B)$, así que $A+B=3$ y $A-B=1$: $A=2$, $B=1$.\par
\textbf{Con valores que anulan factores:} $x=1$ da $4=2A$, $A=2$; $x=-1$ da $-2=-2B$, $B=1$.
\[ \frac{3x+1}{x^2-1}=\frac{2}{x-1}+\frac{1}{x+1}. \]
\end{ejemplo}

\begin{ojo}[Antes de partir]
Si la fracción no es propia ($\gr P\ge\gr Q$), primero se divide: $\frac{P}{Q}=C(x)+\frac{R}{Q}$, y solo se parte
$\frac{R}{Q}$. Y un factor repetido necesita \emph{una fracción por cada potencia}: $\frac{1}{(x-1)^2}$ no se
escribe solo como $\frac{A}{(x-1)^2}$ cuando acompaña a otros factores.
\end{ojo}

%=====================================================================
\clearpage
\repaso
\begin{center}\renewcommand{\arraystretch}{1.55}
\begin{tabularx}{\linewidth}{@{}l L@{}}
\toprule
Idea & Fórmula o regla\\ \midrule
Dominio de $\frac{P}{Q}$ & $\conjunto{x\in\R\tq Q(x)\neq0}$, leído en la expresión original\\
Simplificar & $\frac{AC}{BC}=\frac{A}{B}$ si $C\neq0$; solo se cancelan factores\\
Asíntota vertical & ceros del denominador que no se cancelan; los que se cancelan dan huecos\\
Asíntota horizontal & $n<m$: $y=0$;\quad $n=m$: cociente de coeficientes principales;\quad $n>m$: no hay\\
Producto y cociente & $\frac{A}{B}\cdot\frac{C}{D}=\frac{AC}{BD}$,\quad $\frac{A}{B}\div\frac{C}{D}=\frac{A}{B}\cdot\frac{D}{C}$\\
Suma y resta & mcd: cada factor una vez con su mayor exponente; se amplifica lo que falta\\
Sistema lineal $2\times2$ & corte de dos rectas: una, ninguna o infinitas soluciones\\
Métodos & reducción, sustitución, igualación; escalonar en $3\times3$\\
Sistema no lineal & despejar, sustituir y resolver la ecuación que queda\\
Fracciones parciales & $\frac{A}{x-a}$;\quad $\frac{A_1}{x-a}+\dots+\frac{A_k}{(x-a)^k}$;\quad $\frac{Bx+C}{ax^2+bx+c}$\\
\bottomrule
\end{tabularx}
\end{center}

\begin{ojo}[Errores frecuentes]
\begin{itemize}
\item Cancelar términos que se suman: $\frac{x+5}{5}\neq x$.
\item Dar el dominio de la expresión ya simplificada y olvidar el valor del factor cancelado.
\item Sumar numeradores con numeradores y denominadores con denominadores.
\item Olvidar el paréntesis al restar una fracción cuyo numerador tiene varios términos.
\item No comprobar la solución en todas las ecuaciones del sistema.
\item Partir una fracción impropia sin dividir antes, o poner una constante sobre un factor cuadrático.
\end{itemize}
\end{ojo}

%=====================================================================
\taller
\begin{ejercicios}
\item Halla el dominio y simplifica:
  \begin{ejercicios}
  \item $\dfrac{x^2-5x+6}{x^2-9}$ \qquad (b) $\dfrac{2x^2+6x}{x^2+6x+9}$ \qquad (c) $\dfrac{x^3-8}{x-2}$
  \end{ejercicios}
\item Sea $f(x)=\dfrac{x^2-x-2}{x^2-4}$.
  \begin{ejercicios}
  \item Halla su dominio, sus asíntotas y su hueco (con sus coordenadas).
  \item Dibuja la gráfica de $f$ marcando las asíntotas con líneas punteadas y el hueco con un círculo vacío.
  \end{ejercicios}
\item Efectúa y simplifica:
  \begin{ejercicios}
  \item $\dfrac{x^2-1}{x^2+2x}\cdot\dfrac{x}{x+1}$ \qquad (b) $\dfrac{x^2-4}{3x}\div\dfrac{x+2}{6x^2}$
  \end{ejercicios}
\item Efectúa y simplifica:
  \begin{ejercicios}
  \item $\dfrac{3}{x-2}-\dfrac{2}{x+1}$ \qquad (b) $\dfrac{x}{x^2-1}+\dfrac{1}{x+1}$
  \end{ejercicios}
\item Simplifica la fracción compleja $\dfrac{\frac1x-\frac1y}{\frac1{x^2}-\frac1{y^2}}$.
\item Una llave llena un tanque en $4$ horas y un desagüe lo vacía en $6$ horas. Si el tanque está vacío y se abren
  los dos a la vez, ¿en cuánto tiempo se llena?
\item Para cada sistema, di si tiene una, ninguna o infinitas soluciones y, si tiene una, hállala.
  \begin{ejercicios}
  \item $3x-2y=4,\quad 5x+3y=13$
  \item $2x-4y=6,\quad -x+2y=3$
  \item $6x+9y=12,\quad 2x+3y=4$
  \end{ejercicios}
\item A una función de teatro entraron $100$ personas y se recaudaron $\$900\,000$. La boleta de adulto cuesta
  $\$12\,000$, la de estudiante $\$8\,000$ y la de niño $\$5\,000$. Entraron el doble de estudiantes que de niños.
  Plantea un sistema de tres ecuaciones y halla cuántas boletas de cada tipo se vendieron.
\item Sistemas no lineales.
  \begin{ejercicios}
  \item Resuelve el sistema $x^2+y^2=25$, $\ x+y=7$ e interpreta la respuesta en una gráfica (una circunferencia y una recta).
  \item ¿Para qué valor de $b$ la recta $y=2x+b$ toca la parábola $y=x^2$ en un solo punto?
  \end{ejercicios}
\item Descompón en fracciones parciales:
  \begin{ejercicios}
  \item $\dfrac{5x}{x^2+x-6}$ \qquad (b) $\dfrac{3x+5}{(x+1)^2}$ \qquad (c) $\dfrac{2x^2+x+3}{(x+1)\left(x^2+1\right)}$
  \end{ejercicios}
\end{ejercicios}

\clearpage
\respuestas
\begin{respuestascortas}
\begin{soluciones}
\item (a) $x\neq\pm3$; $\dfrac{x-2}{x+3}$. (b) $x\neq-3$; $\dfrac{2x}{x+3}$. (c) $x\neq2$; $x^2+2x+4$.
\item (a) $f(x)=\dfrac{(x-2)(x+1)}{(x-2)(x+2)}=\dfrac{x+1}{x+2}$, dominio $\R\setminus\{\pm2\}$; asíntotas $x=-2$ y
  $y=1$; hueco en $\left(2,\tfrac34\right)$. (b) Hipérbola con ramas a los lados de $x=-2$; corta los ejes en
  $(-1,0)$ y $\left(0,\tfrac12\right)$; para $x>-2$ queda por debajo de $y=1$ y para $x<-2$, por encima.
\item (a) $\dfrac{x-1}{x+2}$, con $x\neq0,-1,-2$. (b) $2x(x-2)$, con $x\neq0,-2$.
\item (a) $\dfrac{x+7}{(x-2)(x+1)}$. (b) $\dfrac{2x-1}{x^2-1}$.
\item $\dfrac{xy}{x+y}$ (con $x\neq\pm y$, $x,y\neq0$).
\item Por hora: $\frac14-\frac16=\frac1{12}$ de tanque; se llena en $12$ horas.
\item (a) Una: $(2,1)$. (b) Ninguna: $x-2y=3$ y $x-2y=-3$ son paralelas. (c) Infinitas: la primera es el triple de la segunda.
\item $a+e+n=100$, $12a+8e+5n=900$ (en miles), $e=2n$. Solución: $40$ adultos, $40$ estudiantes y $20$ niños.
\item (a) $y=7-x$ da $x^2-7x+12=0$: los puntos $(3,4)$ y $(4,3)$, donde la recta corta la circunferencia de radio $5$.
  (b) $x^2-2x-b=0$ con discriminante $4+4b=0$: $b=-1$ (tangente en $(1,1)$).
\item (a) $\dfrac{2}{x-2}+\dfrac{3}{x+3}$. (b) $\dfrac{3}{x+1}+\dfrac{2}{(x+1)^2}$. (c) $\dfrac{2}{x+1}+\dfrac{1}{x^2+1}$.
\end{soluciones}
\end{respuestascortas}

\end{document}
